% ============================================================================
% chap5_architecture.tex
%   第5章 · 多机器人安全强化学习的容器化仿真验证平台
%   五层架构图 (基础平台 + 本文扩展)
%
% 独立编译:
%   xelatex figs/chap5_architecture.tex
%   → figs/chap5_architecture.pdf
%
% 嵌入论文:
%   \begin{figure}[htbp]
%     \centering
%     \resizebox{\linewidth}{!}{% ============================================================================
% chap5_architecture.tex
%   第5章 · 多机器人安全强化学习的容器化仿真验证平台
%   五层架构图 (基础平台 + 本文扩展)
%
% 独立编译:
%   xelatex figs/chap5_architecture.tex
%   → figs/chap5_architecture.pdf
%
% 嵌入论文:
%   \begin{figure}[htbp]
%     \centering
%     \resizebox{\linewidth}{!}{% ============================================================================
% chap5_architecture.tex
%   第5章 · 多机器人安全强化学习的容器化仿真验证平台
%   五层架构图 (基础平台 + 本文扩展)
%
% 独立编译:
%   xelatex figs/chap5_architecture.tex
%   → figs/chap5_architecture.pdf
%
% 嵌入论文:
%   \begin{figure}[htbp]
%     \centering
%     \resizebox{\linewidth}{!}{% ============================================================================
% chap5_architecture.tex
%   第5章 · 多机器人安全强化学习的容器化仿真验证平台
%   五层架构图 (基础平台 + 本文扩展)
%
% 独立编译:
%   xelatex figs/chap5_architecture.tex
%   → figs/chap5_architecture.pdf
%
% 嵌入论文:
%   \begin{figure}[htbp]
%     \centering
%     \resizebox{\linewidth}{!}{\input{figs/chap5_architecture}}
%     \caption{多机器人安全强化学习容器化仿真验证平台的五层架构。
%              蓝色/橙色/绿色实线方框为本文扩展实现, 灰色虚线方框为通用基础模块.}
%     \label{fig:chap5_architecture}
%   \end{figure}
% ============================================================================

\documentclass[tikz,border=6pt]{standalone}
\usepackage[UTF8]{ctex}
\usepackage{tikz}
\usetikzlibrary{
  positioning,
  arrows.meta,
  shapes.geometric,
  shapes.multipart,
  calc,
  backgrounds,
  fit,
  decorations.pathmorphing,
}

% 颜色整体加深：边框与箭头颜色保持不变，文字统一改为黑色
\definecolor{cOurs}{HTML}{0B4F8A}      % 深蓝：本文核心扩展
\definecolor{cOursAlt}{HTML}{C05600}   % 深橙：本文方法集成
\definecolor{cOursEval}{HTML}{000000}  % 黑色：本文评估管线
\definecolor{cUpstream}{HTML}{4D4D4D}  % 深灰：通用部件
\definecolor{cBgLayer}{HTML}{F0F0F0}
\definecolor{cBgHighlight}{HTML}{FFF3E0}

\begin{document}

\begin{tikzpicture}[
  font=\large\sffamily,
  every node/.style={align=center, text=black},
  % 层容器
  layer/.style={
    draw=black!40, thick, rounded corners=2pt,
    fill=cBgLayer, inner sep=8pt,
    minimum width=18cm, minimum height=1.75cm,
  },
  layerLabel/.style={
    font=\bfseries\large, text=black, align=left, anchor=west,
  },
  % 模块 (本文扩展)
  ours/.style={
    draw=cOurs, line width=1pt, rounded corners=2pt,
    fill=white, text=black, inner sep=4pt,
    minimum width=3.5cm, minimum height=0.9cm,
    font=\normalsize\sffamily,
  },
  oursAlt/.style={
    draw=cOursAlt, line width=1pt, rounded corners=2pt,
    fill=white, text=black, inner sep=4pt,
    minimum width=3.5cm, minimum height=0.9cm,
    font=\normalsize\sffamily,
  },
  oursEval/.style={
    draw=cOursEval, line width=1pt, rounded corners=2pt,
    fill=white, text=black, inner sep=4pt,
    minimum width=3.5cm, minimum height=0.9cm,
    font=\normalsize\sffamily,
  },
  % 模块 (通用部件)
  upstream/.style={
    draw=cUpstream, line width=0.6pt, dashed, rounded corners=2pt,
    fill=white, text=black, inner sep=4pt,
    minimum width=3.5cm, minimum height=0.9cm,
    font=\normalsize\sffamily,
  },
  % 流向箭头
  flow/.style={
    -{Stealth[length=2.2mm, width=1.8mm]},
    line width=0.9pt, draw=black!75,
  },
  flowBold/.style={
    -{Stealth[length=2.6mm, width=2.2mm]},
    line width=1.3pt, draw=cOursAlt,
  },
  codepath/.style={font=\footnotesize\ttfamily, text=black},
]

% ============================================================================
% 五层, 自底向上排列
% ============================================================================

% 4 列 x 坐标(相对层 center), 统一用同一套值保证上下列对齐
\newcommand{\colA}{-6.3cm}
\newcommand{\colB}{-2.1cm}
\newcommand{\colC}{2.1cm}
\newcommand{\colD}{6.3cm}

% ---- Layer 1: 仿真场景层 ----
\node[layer] (L1) at (0, 0) {};
\node[layerLabel] at ($(L1.north west) + (0.15cm, -0.25cm)$) {Layer 1 · 仿真场景层 (Webots)};
\node[ours] (L1A) at ($(L1.center) + (\colA, -0.20cm)$) {E-puck 差速模型};
\node[ours] (L1B) at ($(L1.center) + (\colB, -0.20cm)$) {多机场景};
\node[ours] (L1C) at ($(L1.center) + (\colC, -0.20cm)$) {可配置编队};
\node[ours] (L1D) at ($(L1.center) + (\colD, -0.20cm)$) {静态障碍 + 目标};

% ---- Layer 2: 控制接口层 (ROS2) ----
\node[layer, above=0.55cm of L1] (L2) {};
\node[layerLabel] at ($(L2.north west) + (0.15cm, -0.25cm)$) {Layer 2 · 控制接口层 (ROS2)};
\node[upstream] (L2A) at ($(L2.center) + (\colA, -0.20cm)$) {ROS2 Humble};
\node[ours] (L2B) at ($(L2.center) + (\colB, -0.20cm)$) {多机话题契约};
\node[ours] (L2C) at ($(L2.center) + (\colC, -0.20cm)$) {指令分发};
\node[ours] (L2D) at ($(L2.center) + (\colD, -0.20cm)$) {状态广播};

% ---- Layer 3: 算法容器层 (Docker) ----
\node[layer, above=0.55cm of L2] (L3) {};
\node[layerLabel] at ($(L3.north west) + (0.15cm, -0.25cm)$) {Layer 3 · 算法容器层 (Docker)};
\node[upstream] (L3A) at ($(L3.center) + (\colA, -0.20cm)$) {Docker + noVNC};
\node[ours] (L3B) at ($(L3.center) + (\colB, -0.20cm)$) {PyTorch + CUDA};
\node[ours] (L3C) at ($(L3.center) + (\colC, -0.20cm)$) {safepo + safety-gym};
\node[ours] (L3D) at ($(L3.center) + (\colD, -0.20cm)$) {MAPPO 检查点共享};

% ---- Layer 4: 分层集成层 (核心) ----
\node[layer, minimum height=2.6cm, above=0.55cm of L3,
      fill=cBgHighlight, draw=cOursAlt, line width=1pt] (L4) {};
\node[layerLabel, text=black]
  at ($(L4.north west) + (0.15cm, -0.25cm)$)
  {Layer 4 · 分层控制栈集成层 \emph{(本文方法贡献)}};

% 四个集成模块 - 用统一列坐标, 避免重叠
\node[oursAlt] (L4A) at ($(L4.center) + (\colA, -0.35cm)$)
  {MAPPO 策略加载};
\node[oursAlt] (L4B) at ($(L4.center) + (\colB, -0.35cm)$)
  {GCPL 软协调层};
\node[oursAlt] (L4C) at ($(L4.center) + (\colC, -0.35cm)$)
  {CMPL 硬投影层};
\node[oursAlt] (L4D) at ($(L4.center) + (\colD, -0.35cm)$)
  {差速映射};

% Supervisor 容器
\node[draw=cOursAlt, dash pattern=on 2pt off 1pt, rounded corners=3pt,
      fit=(L4A) (L4B) (L4C) (L4D), inner sep=5pt,
      label={[text=black, font=\small\bfseries, fill=white, inner sep=1.5pt, yshift=-3pt]
             below:gcpl\_full\_stack\_supervisor.py}] (L4box) {};

% Layer 4 内部的流向箭头
\draw[flowBold] (L4A.east) -- node[above, font=\footnotesize, text=black, fill=white, inner sep=1pt]{$a^{RL}$} (L4B.west);
\draw[flowBold] (L4B.east) -- node[above, font=\footnotesize, text=black, fill=white, inner sep=1pt]{$a^\star$} (L4C.west);
\draw[flowBold] (L4C.east) -- node[above, font=\footnotesize, text=black, fill=white, inner sep=1pt]{$a_\mathrm{safe}$} (L4D.west);

% ---- Layer 5: 评估管线层 ----
\node[layer, above=0.55cm of L4] (L5) {};
\node[layerLabel] at ($(L5.north west) + (0.15cm, -0.25cm)$) {Layer 5 · 批量评估与可视化管线};
\node[oursEval] (L5A) at ($(L5.center) + (\colA, -0.20cm)$)
  {对比 sweep};
\node[oursEval] (L5B) at ($(L5.center) + (\colB, -0.20cm)$)
  {消融 sweep};
\node[oursEval] (L5C) at ($(L5.center) + (\colC, -0.20cm)$)
  {批量评估};
\node[oursEval] (L5D) at ($(L5.center) + (\colD, -0.20cm)$)
  {多 seed 绘图};

% ============================================================================
% 跨层数据流
% ============================================================================

% ----  跨层主链: 走 layer 框外侧, 用平滑曲线 (Bezier) 贯穿 L1↔L4  ----
%        左侧上行 L1 → L4 (状态), 右侧下行 L4 → L1 (指令)
\draw[flow, rounded corners=8pt]
  ($(L1.west)+(-0.6cm,0)$)
    .. controls ($(L1.west)+(-1.6cm,1.0cm)$) and ($(L4.west)+(-1.6cm,-1.0cm)$) ..
  node[right, pos=0.5, xshift=2pt, font=\footnotesize, text=black, fill=white, inner sep=1pt]{状态上行}
  ($(L4.west)+(-0.6cm,0)$);

\draw[flow, rounded corners=8pt]
  ($(L4.east)+(0.6cm,0)$)
    .. controls ($(L4.east)+(1.6cm,-1.0cm)$) and ($(L1.east)+(1.6cm,1.0cm)$) ..
  node[left, pos=0.5, xshift=-2pt, font=\footnotesize, text=black, fill=white, inner sep=1pt]{指令下行}
  ($(L1.east)+(0.6cm,0)$);

% ----  L3 → L4  (容器提供 runtime, 虚线纵向, 列对齐)  ----
\draw[flow, dashed] (L3A.north) -- (L4A.south);
\draw[flow, dashed] (L3B.north) -- (L4B.south);
\draw[flow, dashed] (L3C.north) -- (L4C.south);
\draw[flow, dashed] (L3D.north) -- (L4D.south);

% ----  L5 → L4  (评估管线触发, 虚线纵向, 列对齐)  ----
\draw[flow, dashed] (L5A.south) -- (L4A.north);
\draw[flow, dashed] (L5B.south) -- (L4B.north);
\draw[flow, dashed] (L5C.south) -- (L4C.north);
\draw[flow, dashed] (L5D.south) -- (L4D.north);

% ============================================================================
% Legend  (放在图下方, 横向排布, 避开五层主体)
% ============================================================================
\node[draw=black!30, rounded corners=2pt, inner sep=5pt, font=\small, text=black,
      anchor=north]
  at ($(L1.south) + (0, -0.4cm)$) (legend)
  {\begin{tabular}{@{}l@{\hspace{3pt}}l@{\hspace{10pt}}l@{\hspace{3pt}}l@{\hspace{10pt}}l@{\hspace{3pt}}l@{\hspace{10pt}}l@{\hspace{3pt}}l@{\hspace{10pt}}l@{\hspace{3pt}}l@{\hspace{10pt}}l@{\hspace{3pt}}l@{\hspace{10pt}}l@{\hspace{3pt}}l@{}}
    \tikz \draw[cOursAlt, line width=1pt, fill=white, rounded corners=1pt]
        (0,0) rectangle (0.4,0.2); & 本文方法贡献 (L4) &
    \tikz \draw[cOurs, line width=1pt, fill=white, rounded corners=1pt]
        (0,0) rectangle (0.4,0.2); & 本文平台扩展 (L1--L3) &
    \tikz \draw[cOursEval, line width=1pt, fill=white, rounded corners=1pt]
        (0,0) rectangle (0.4,0.2); & 评估管线扩展 (L5) &
    \tikz \draw[cUpstream, dashed, fill=white, rounded corners=1pt]
        (0,0) rectangle (0.4,0.2); & 通用基础模块 &
    \tikz \draw[flowBold, shorten >=2pt] (0,0.1) -- (0.4,0.1); & 核心控制管道 &
    \tikz \draw[flow, shorten >=2pt] (0,0.1) -- (0.4,0.1); & 数据/状态流向 &
    \tikz \draw[flow, dashed, shorten >=2pt] (0,0.1) -- (0.4,0.1); & 依赖关系 \\
  \end{tabular}};

\end{tikzpicture}

\end{document}}
%     \caption{多机器人安全强化学习容器化仿真验证平台的五层架构。
%              蓝色/橙色/绿色实线方框为本文扩展实现, 灰色虚线方框为通用基础模块.}
%     \label{fig:chap5_architecture}
%   \end{figure}
% ============================================================================

\documentclass[tikz,border=6pt]{standalone}
\usepackage[UTF8]{ctex}
\usepackage{tikz}
\usetikzlibrary{
  positioning,
  arrows.meta,
  shapes.geometric,
  shapes.multipart,
  calc,
  backgrounds,
  fit,
  decorations.pathmorphing,
}

% 颜色整体加深：边框与箭头颜色保持不变，文字统一改为黑色
\definecolor{cOurs}{HTML}{0B4F8A}      % 深蓝：本文核心扩展
\definecolor{cOursAlt}{HTML}{C05600}   % 深橙：本文方法集成
\definecolor{cOursEval}{HTML}{000000}  % 黑色：本文评估管线
\definecolor{cUpstream}{HTML}{4D4D4D}  % 深灰：通用部件
\definecolor{cBgLayer}{HTML}{F0F0F0}
\definecolor{cBgHighlight}{HTML}{FFF3E0}

\begin{document}

\begin{tikzpicture}[
  font=\large\sffamily,
  every node/.style={align=center, text=black},
  % 层容器
  layer/.style={
    draw=black!40, thick, rounded corners=2pt,
    fill=cBgLayer, inner sep=8pt,
    minimum width=18cm, minimum height=1.75cm,
  },
  layerLabel/.style={
    font=\bfseries\large, text=black, align=left, anchor=west,
  },
  % 模块 (本文扩展)
  ours/.style={
    draw=cOurs, line width=1pt, rounded corners=2pt,
    fill=white, text=black, inner sep=4pt,
    minimum width=3.5cm, minimum height=0.9cm,
    font=\normalsize\sffamily,
  },
  oursAlt/.style={
    draw=cOursAlt, line width=1pt, rounded corners=2pt,
    fill=white, text=black, inner sep=4pt,
    minimum width=3.5cm, minimum height=0.9cm,
    font=\normalsize\sffamily,
  },
  oursEval/.style={
    draw=cOursEval, line width=1pt, rounded corners=2pt,
    fill=white, text=black, inner sep=4pt,
    minimum width=3.5cm, minimum height=0.9cm,
    font=\normalsize\sffamily,
  },
  % 模块 (通用部件)
  upstream/.style={
    draw=cUpstream, line width=0.6pt, dashed, rounded corners=2pt,
    fill=white, text=black, inner sep=4pt,
    minimum width=3.5cm, minimum height=0.9cm,
    font=\normalsize\sffamily,
  },
  % 流向箭头
  flow/.style={
    -{Stealth[length=2.2mm, width=1.8mm]},
    line width=0.9pt, draw=black!75,
  },
  flowBold/.style={
    -{Stealth[length=2.6mm, width=2.2mm]},
    line width=1.3pt, draw=cOursAlt,
  },
  codepath/.style={font=\footnotesize\ttfamily, text=black},
]

% ============================================================================
% 五层, 自底向上排列
% ============================================================================

% 4 列 x 坐标(相对层 center), 统一用同一套值保证上下列对齐
\newcommand{\colA}{-6.3cm}
\newcommand{\colB}{-2.1cm}
\newcommand{\colC}{2.1cm}
\newcommand{\colD}{6.3cm}

% ---- Layer 1: 仿真场景层 ----
\node[layer] (L1) at (0, 0) {};
\node[layerLabel] at ($(L1.north west) + (0.15cm, -0.25cm)$) {Layer 1 · 仿真场景层 (Webots)};
\node[ours] (L1A) at ($(L1.center) + (\colA, -0.20cm)$) {E-puck 差速模型};
\node[ours] (L1B) at ($(L1.center) + (\colB, -0.20cm)$) {多机场景};
\node[ours] (L1C) at ($(L1.center) + (\colC, -0.20cm)$) {可配置编队};
\node[ours] (L1D) at ($(L1.center) + (\colD, -0.20cm)$) {静态障碍 + 目标};

% ---- Layer 2: 控制接口层 (ROS2) ----
\node[layer, above=0.55cm of L1] (L2) {};
\node[layerLabel] at ($(L2.north west) + (0.15cm, -0.25cm)$) {Layer 2 · 控制接口层 (ROS2)};
\node[upstream] (L2A) at ($(L2.center) + (\colA, -0.20cm)$) {ROS2 Humble};
\node[ours] (L2B) at ($(L2.center) + (\colB, -0.20cm)$) {多机话题契约};
\node[ours] (L2C) at ($(L2.center) + (\colC, -0.20cm)$) {指令分发};
\node[ours] (L2D) at ($(L2.center) + (\colD, -0.20cm)$) {状态广播};

% ---- Layer 3: 算法容器层 (Docker) ----
\node[layer, above=0.55cm of L2] (L3) {};
\node[layerLabel] at ($(L3.north west) + (0.15cm, -0.25cm)$) {Layer 3 · 算法容器层 (Docker)};
\node[upstream] (L3A) at ($(L3.center) + (\colA, -0.20cm)$) {Docker + noVNC};
\node[ours] (L3B) at ($(L3.center) + (\colB, -0.20cm)$) {PyTorch + CUDA};
\node[ours] (L3C) at ($(L3.center) + (\colC, -0.20cm)$) {safepo + safety-gym};
\node[ours] (L3D) at ($(L3.center) + (\colD, -0.20cm)$) {MAPPO 检查点共享};

% ---- Layer 4: 分层集成层 (核心) ----
\node[layer, minimum height=2.6cm, above=0.55cm of L3,
      fill=cBgHighlight, draw=cOursAlt, line width=1pt] (L4) {};
\node[layerLabel, text=black]
  at ($(L4.north west) + (0.15cm, -0.25cm)$)
  {Layer 4 · 分层控制栈集成层 \emph{(本文方法贡献)}};

% 四个集成模块 - 用统一列坐标, 避免重叠
\node[oursAlt] (L4A) at ($(L4.center) + (\colA, -0.35cm)$)
  {MAPPO 策略加载};
\node[oursAlt] (L4B) at ($(L4.center) + (\colB, -0.35cm)$)
  {GCPL 软协调层};
\node[oursAlt] (L4C) at ($(L4.center) + (\colC, -0.35cm)$)
  {CMPL 硬投影层};
\node[oursAlt] (L4D) at ($(L4.center) + (\colD, -0.35cm)$)
  {差速映射};

% Supervisor 容器
\node[draw=cOursAlt, dash pattern=on 2pt off 1pt, rounded corners=3pt,
      fit=(L4A) (L4B) (L4C) (L4D), inner sep=5pt,
      label={[text=black, font=\small\bfseries, fill=white, inner sep=1.5pt, yshift=-3pt]
             below:gcpl\_full\_stack\_supervisor.py}] (L4box) {};

% Layer 4 内部的流向箭头
\draw[flowBold] (L4A.east) -- node[above, font=\footnotesize, text=black, fill=white, inner sep=1pt]{$a^{RL}$} (L4B.west);
\draw[flowBold] (L4B.east) -- node[above, font=\footnotesize, text=black, fill=white, inner sep=1pt]{$a^\star$} (L4C.west);
\draw[flowBold] (L4C.east) -- node[above, font=\footnotesize, text=black, fill=white, inner sep=1pt]{$a_\mathrm{safe}$} (L4D.west);

% ---- Layer 5: 评估管线层 ----
\node[layer, above=0.55cm of L4] (L5) {};
\node[layerLabel] at ($(L5.north west) + (0.15cm, -0.25cm)$) {Layer 5 · 批量评估与可视化管线};
\node[oursEval] (L5A) at ($(L5.center) + (\colA, -0.20cm)$)
  {对比 sweep};
\node[oursEval] (L5B) at ($(L5.center) + (\colB, -0.20cm)$)
  {消融 sweep};
\node[oursEval] (L5C) at ($(L5.center) + (\colC, -0.20cm)$)
  {批量评估};
\node[oursEval] (L5D) at ($(L5.center) + (\colD, -0.20cm)$)
  {多 seed 绘图};

% ============================================================================
% 跨层数据流
% ============================================================================

% ----  跨层主链: 走 layer 框外侧, 用平滑曲线 (Bezier) 贯穿 L1↔L4  ----
%        左侧上行 L1 → L4 (状态), 右侧下行 L4 → L1 (指令)
\draw[flow, rounded corners=8pt]
  ($(L1.west)+(-0.6cm,0)$)
    .. controls ($(L1.west)+(-1.6cm,1.0cm)$) and ($(L4.west)+(-1.6cm,-1.0cm)$) ..
  node[right, pos=0.5, xshift=2pt, font=\footnotesize, text=black, fill=white, inner sep=1pt]{状态上行}
  ($(L4.west)+(-0.6cm,0)$);

\draw[flow, rounded corners=8pt]
  ($(L4.east)+(0.6cm,0)$)
    .. controls ($(L4.east)+(1.6cm,-1.0cm)$) and ($(L1.east)+(1.6cm,1.0cm)$) ..
  node[left, pos=0.5, xshift=-2pt, font=\footnotesize, text=black, fill=white, inner sep=1pt]{指令下行}
  ($(L1.east)+(0.6cm,0)$);

% ----  L3 → L4  (容器提供 runtime, 虚线纵向, 列对齐)  ----
\draw[flow, dashed] (L3A.north) -- (L4A.south);
\draw[flow, dashed] (L3B.north) -- (L4B.south);
\draw[flow, dashed] (L3C.north) -- (L4C.south);
\draw[flow, dashed] (L3D.north) -- (L4D.south);

% ----  L5 → L4  (评估管线触发, 虚线纵向, 列对齐)  ----
\draw[flow, dashed] (L5A.south) -- (L4A.north);
\draw[flow, dashed] (L5B.south) -- (L4B.north);
\draw[flow, dashed] (L5C.south) -- (L4C.north);
\draw[flow, dashed] (L5D.south) -- (L4D.north);

% ============================================================================
% Legend  (放在图下方, 横向排布, 避开五层主体)
% ============================================================================
\node[draw=black!30, rounded corners=2pt, inner sep=5pt, font=\small, text=black,
      anchor=north]
  at ($(L1.south) + (0, -0.4cm)$) (legend)
  {\begin{tabular}{@{}l@{\hspace{3pt}}l@{\hspace{10pt}}l@{\hspace{3pt}}l@{\hspace{10pt}}l@{\hspace{3pt}}l@{\hspace{10pt}}l@{\hspace{3pt}}l@{\hspace{10pt}}l@{\hspace{3pt}}l@{\hspace{10pt}}l@{\hspace{3pt}}l@{\hspace{10pt}}l@{\hspace{3pt}}l@{}}
    \tikz \draw[cOursAlt, line width=1pt, fill=white, rounded corners=1pt]
        (0,0) rectangle (0.4,0.2); & 本文方法贡献 (L4) &
    \tikz \draw[cOurs, line width=1pt, fill=white, rounded corners=1pt]
        (0,0) rectangle (0.4,0.2); & 本文平台扩展 (L1--L3) &
    \tikz \draw[cOursEval, line width=1pt, fill=white, rounded corners=1pt]
        (0,0) rectangle (0.4,0.2); & 评估管线扩展 (L5) &
    \tikz \draw[cUpstream, dashed, fill=white, rounded corners=1pt]
        (0,0) rectangle (0.4,0.2); & 通用基础模块 &
    \tikz \draw[flowBold, shorten >=2pt] (0,0.1) -- (0.4,0.1); & 核心控制管道 &
    \tikz \draw[flow, shorten >=2pt] (0,0.1) -- (0.4,0.1); & 数据/状态流向 &
    \tikz \draw[flow, dashed, shorten >=2pt] (0,0.1) -- (0.4,0.1); & 依赖关系 \\
  \end{tabular}};

\end{tikzpicture}

\end{document}}
%     \caption{多机器人安全强化学习容器化仿真验证平台的五层架构。
%              蓝色/橙色/绿色实线方框为本文扩展实现, 灰色虚线方框为通用基础模块.}
%     \label{fig:chap5_architecture}
%   \end{figure}
% ============================================================================

\documentclass[tikz,border=6pt]{standalone}
\usepackage[UTF8]{ctex}
\usepackage{tikz}
\usetikzlibrary{
  positioning,
  arrows.meta,
  shapes.geometric,
  shapes.multipart,
  calc,
  backgrounds,
  fit,
  decorations.pathmorphing,
}

% 颜色整体加深：边框与箭头颜色保持不变，文字统一改为黑色
\definecolor{cOurs}{HTML}{0B4F8A}      % 深蓝：本文核心扩展
\definecolor{cOursAlt}{HTML}{C05600}   % 深橙：本文方法集成
\definecolor{cOursEval}{HTML}{000000}  % 黑色：本文评估管线
\definecolor{cUpstream}{HTML}{4D4D4D}  % 深灰：通用部件
\definecolor{cBgLayer}{HTML}{F0F0F0}
\definecolor{cBgHighlight}{HTML}{FFF3E0}

\begin{document}

\begin{tikzpicture}[
  font=\large\sffamily,
  every node/.style={align=center, text=black},
  % 层容器
  layer/.style={
    draw=black!40, thick, rounded corners=2pt,
    fill=cBgLayer, inner sep=8pt,
    minimum width=18cm, minimum height=1.75cm,
  },
  layerLabel/.style={
    font=\bfseries\large, text=black, align=left, anchor=west,
  },
  % 模块 (本文扩展)
  ours/.style={
    draw=cOurs, line width=1pt, rounded corners=2pt,
    fill=white, text=black, inner sep=4pt,
    minimum width=3.5cm, minimum height=0.9cm,
    font=\normalsize\sffamily,
  },
  oursAlt/.style={
    draw=cOursAlt, line width=1pt, rounded corners=2pt,
    fill=white, text=black, inner sep=4pt,
    minimum width=3.5cm, minimum height=0.9cm,
    font=\normalsize\sffamily,
  },
  oursEval/.style={
    draw=cOursEval, line width=1pt, rounded corners=2pt,
    fill=white, text=black, inner sep=4pt,
    minimum width=3.5cm, minimum height=0.9cm,
    font=\normalsize\sffamily,
  },
  % 模块 (通用部件)
  upstream/.style={
    draw=cUpstream, line width=0.6pt, dashed, rounded corners=2pt,
    fill=white, text=black, inner sep=4pt,
    minimum width=3.5cm, minimum height=0.9cm,
    font=\normalsize\sffamily,
  },
  % 流向箭头
  flow/.style={
    -{Stealth[length=2.2mm, width=1.8mm]},
    line width=0.9pt, draw=black!75,
  },
  flowBold/.style={
    -{Stealth[length=2.6mm, width=2.2mm]},
    line width=1.3pt, draw=cOursAlt,
  },
  codepath/.style={font=\footnotesize\ttfamily, text=black},
]

% ============================================================================
% 五层, 自底向上排列
% ============================================================================

% 4 列 x 坐标(相对层 center), 统一用同一套值保证上下列对齐
\newcommand{\colA}{-6.3cm}
\newcommand{\colB}{-2.1cm}
\newcommand{\colC}{2.1cm}
\newcommand{\colD}{6.3cm}

% ---- Layer 1: 仿真场景层 ----
\node[layer] (L1) at (0, 0) {};
\node[layerLabel] at ($(L1.north west) + (0.15cm, -0.25cm)$) {Layer 1 · 仿真场景层 (Webots)};
\node[ours] (L1A) at ($(L1.center) + (\colA, -0.20cm)$) {E-puck 差速模型};
\node[ours] (L1B) at ($(L1.center) + (\colB, -0.20cm)$) {多机场景};
\node[ours] (L1C) at ($(L1.center) + (\colC, -0.20cm)$) {可配置编队};
\node[ours] (L1D) at ($(L1.center) + (\colD, -0.20cm)$) {静态障碍 + 目标};

% ---- Layer 2: 控制接口层 (ROS2) ----
\node[layer, above=0.55cm of L1] (L2) {};
\node[layerLabel] at ($(L2.north west) + (0.15cm, -0.25cm)$) {Layer 2 · 控制接口层 (ROS2)};
\node[upstream] (L2A) at ($(L2.center) + (\colA, -0.20cm)$) {ROS2 Humble};
\node[ours] (L2B) at ($(L2.center) + (\colB, -0.20cm)$) {多机话题契约};
\node[ours] (L2C) at ($(L2.center) + (\colC, -0.20cm)$) {指令分发};
\node[ours] (L2D) at ($(L2.center) + (\colD, -0.20cm)$) {状态广播};

% ---- Layer 3: 算法容器层 (Docker) ----
\node[layer, above=0.55cm of L2] (L3) {};
\node[layerLabel] at ($(L3.north west) + (0.15cm, -0.25cm)$) {Layer 3 · 算法容器层 (Docker)};
\node[upstream] (L3A) at ($(L3.center) + (\colA, -0.20cm)$) {Docker + noVNC};
\node[ours] (L3B) at ($(L3.center) + (\colB, -0.20cm)$) {PyTorch + CUDA};
\node[ours] (L3C) at ($(L3.center) + (\colC, -0.20cm)$) {safepo + safety-gym};
\node[ours] (L3D) at ($(L3.center) + (\colD, -0.20cm)$) {MAPPO 检查点共享};

% ---- Layer 4: 分层集成层 (核心) ----
\node[layer, minimum height=2.6cm, above=0.55cm of L3,
      fill=cBgHighlight, draw=cOursAlt, line width=1pt] (L4) {};
\node[layerLabel, text=black]
  at ($(L4.north west) + (0.15cm, -0.25cm)$)
  {Layer 4 · 分层控制栈集成层 \emph{(本文方法贡献)}};

% 四个集成模块 - 用统一列坐标, 避免重叠
\node[oursAlt] (L4A) at ($(L4.center) + (\colA, -0.35cm)$)
  {MAPPO 策略加载};
\node[oursAlt] (L4B) at ($(L4.center) + (\colB, -0.35cm)$)
  {GCPL 软协调层};
\node[oursAlt] (L4C) at ($(L4.center) + (\colC, -0.35cm)$)
  {CMPL 硬投影层};
\node[oursAlt] (L4D) at ($(L4.center) + (\colD, -0.35cm)$)
  {差速映射};

% Supervisor 容器
\node[draw=cOursAlt, dash pattern=on 2pt off 1pt, rounded corners=3pt,
      fit=(L4A) (L4B) (L4C) (L4D), inner sep=5pt,
      label={[text=black, font=\small\bfseries, fill=white, inner sep=1.5pt, yshift=-3pt]
             below:gcpl\_full\_stack\_supervisor.py}] (L4box) {};

% Layer 4 内部的流向箭头
\draw[flowBold] (L4A.east) -- node[above, font=\footnotesize, text=black, fill=white, inner sep=1pt]{$a^{RL}$} (L4B.west);
\draw[flowBold] (L4B.east) -- node[above, font=\footnotesize, text=black, fill=white, inner sep=1pt]{$a^\star$} (L4C.west);
\draw[flowBold] (L4C.east) -- node[above, font=\footnotesize, text=black, fill=white, inner sep=1pt]{$a_\mathrm{safe}$} (L4D.west);

% ---- Layer 5: 评估管线层 ----
\node[layer, above=0.55cm of L4] (L5) {};
\node[layerLabel] at ($(L5.north west) + (0.15cm, -0.25cm)$) {Layer 5 · 批量评估与可视化管线};
\node[oursEval] (L5A) at ($(L5.center) + (\colA, -0.20cm)$)
  {对比 sweep};
\node[oursEval] (L5B) at ($(L5.center) + (\colB, -0.20cm)$)
  {消融 sweep};
\node[oursEval] (L5C) at ($(L5.center) + (\colC, -0.20cm)$)
  {批量评估};
\node[oursEval] (L5D) at ($(L5.center) + (\colD, -0.20cm)$)
  {多 seed 绘图};

% ============================================================================
% 跨层数据流
% ============================================================================

% ----  跨层主链: 走 layer 框外侧, 用平滑曲线 (Bezier) 贯穿 L1↔L4  ----
%        左侧上行 L1 → L4 (状态), 右侧下行 L4 → L1 (指令)
\draw[flow, rounded corners=8pt]
  ($(L1.west)+(-0.6cm,0)$)
    .. controls ($(L1.west)+(-1.6cm,1.0cm)$) and ($(L4.west)+(-1.6cm,-1.0cm)$) ..
  node[right, pos=0.5, xshift=2pt, font=\footnotesize, text=black, fill=white, inner sep=1pt]{状态上行}
  ($(L4.west)+(-0.6cm,0)$);

\draw[flow, rounded corners=8pt]
  ($(L4.east)+(0.6cm,0)$)
    .. controls ($(L4.east)+(1.6cm,-1.0cm)$) and ($(L1.east)+(1.6cm,1.0cm)$) ..
  node[left, pos=0.5, xshift=-2pt, font=\footnotesize, text=black, fill=white, inner sep=1pt]{指令下行}
  ($(L1.east)+(0.6cm,0)$);

% ----  L3 → L4  (容器提供 runtime, 虚线纵向, 列对齐)  ----
\draw[flow, dashed] (L3A.north) -- (L4A.south);
\draw[flow, dashed] (L3B.north) -- (L4B.south);
\draw[flow, dashed] (L3C.north) -- (L4C.south);
\draw[flow, dashed] (L3D.north) -- (L4D.south);

% ----  L5 → L4  (评估管线触发, 虚线纵向, 列对齐)  ----
\draw[flow, dashed] (L5A.south) -- (L4A.north);
\draw[flow, dashed] (L5B.south) -- (L4B.north);
\draw[flow, dashed] (L5C.south) -- (L4C.north);
\draw[flow, dashed] (L5D.south) -- (L4D.north);

% ============================================================================
% Legend  (放在图下方, 横向排布, 避开五层主体)
% ============================================================================
\node[draw=black!30, rounded corners=2pt, inner sep=5pt, font=\small, text=black,
      anchor=north]
  at ($(L1.south) + (0, -0.4cm)$) (legend)
  {\begin{tabular}{@{}l@{\hspace{3pt}}l@{\hspace{10pt}}l@{\hspace{3pt}}l@{\hspace{10pt}}l@{\hspace{3pt}}l@{\hspace{10pt}}l@{\hspace{3pt}}l@{\hspace{10pt}}l@{\hspace{3pt}}l@{\hspace{10pt}}l@{\hspace{3pt}}l@{\hspace{10pt}}l@{\hspace{3pt}}l@{}}
    \tikz \draw[cOursAlt, line width=1pt, fill=white, rounded corners=1pt]
        (0,0) rectangle (0.4,0.2); & 本文方法贡献 (L4) &
    \tikz \draw[cOurs, line width=1pt, fill=white, rounded corners=1pt]
        (0,0) rectangle (0.4,0.2); & 本文平台扩展 (L1--L3) &
    \tikz \draw[cOursEval, line width=1pt, fill=white, rounded corners=1pt]
        (0,0) rectangle (0.4,0.2); & 评估管线扩展 (L5) &
    \tikz \draw[cUpstream, dashed, fill=white, rounded corners=1pt]
        (0,0) rectangle (0.4,0.2); & 通用基础模块 &
    \tikz \draw[flowBold, shorten >=2pt] (0,0.1) -- (0.4,0.1); & 核心控制管道 &
    \tikz \draw[flow, shorten >=2pt] (0,0.1) -- (0.4,0.1); & 数据/状态流向 &
    \tikz \draw[flow, dashed, shorten >=2pt] (0,0.1) -- (0.4,0.1); & 依赖关系 \\
  \end{tabular}};

\end{tikzpicture}

\end{document}}
%     \caption{多机器人安全强化学习容器化仿真验证平台的五层架构。
%              蓝色/橙色/绿色实线方框为本文扩展实现, 灰色虚线方框为通用基础模块.}
%     \label{fig:chap5_architecture}
%   \end{figure}
% ============================================================================

\documentclass[tikz,border=6pt]{standalone}
\usepackage[UTF8]{ctex}
\usepackage{tikz}
\usetikzlibrary{
  positioning,
  arrows.meta,
  shapes.geometric,
  shapes.multipart,
  calc,
  backgrounds,
  fit,
  decorations.pathmorphing,
}

% 颜色整体加深：边框与箭头颜色保持不变，文字统一改为黑色
\definecolor{cOurs}{HTML}{0B4F8A}      % 深蓝：本文核心扩展
\definecolor{cOursAlt}{HTML}{C05600}   % 深橙：本文方法集成
\definecolor{cOursEval}{HTML}{000000}  % 黑色：本文评估管线
\definecolor{cUpstream}{HTML}{4D4D4D}  % 深灰：通用部件
\definecolor{cBgLayer}{HTML}{F0F0F0}
\definecolor{cBgHighlight}{HTML}{FFF3E0}

\begin{document}

\begin{tikzpicture}[
  font=\large\sffamily,
  every node/.style={align=center, text=black},
  % 层容器
  layer/.style={
    draw=black!40, thick, rounded corners=2pt,
    fill=cBgLayer, inner sep=8pt,
    minimum width=18cm, minimum height=1.75cm,
  },
  layerLabel/.style={
    font=\bfseries\large, text=black, align=left, anchor=west,
  },
  % 模块 (本文扩展)
  ours/.style={
    draw=cOurs, line width=1pt, rounded corners=2pt,
    fill=white, text=black, inner sep=4pt,
    minimum width=3.5cm, minimum height=0.9cm,
    font=\normalsize\sffamily,
  },
  oursAlt/.style={
    draw=cOursAlt, line width=1pt, rounded corners=2pt,
    fill=white, text=black, inner sep=4pt,
    minimum width=3.5cm, minimum height=0.9cm,
    font=\normalsize\sffamily,
  },
  oursEval/.style={
    draw=cOursEval, line width=1pt, rounded corners=2pt,
    fill=white, text=black, inner sep=4pt,
    minimum width=3.5cm, minimum height=0.9cm,
    font=\normalsize\sffamily,
  },
  % 模块 (通用部件)
  upstream/.style={
    draw=cUpstream, line width=0.6pt, dashed, rounded corners=2pt,
    fill=white, text=black, inner sep=4pt,
    minimum width=3.5cm, minimum height=0.9cm,
    font=\normalsize\sffamily,
  },
  % 流向箭头
  flow/.style={
    -{Stealth[length=2.2mm, width=1.8mm]},
    line width=0.9pt, draw=black!75,
  },
  flowBold/.style={
    -{Stealth[length=2.6mm, width=2.2mm]},
    line width=1.3pt, draw=cOursAlt,
  },
  codepath/.style={font=\footnotesize\ttfamily, text=black},
]

% ============================================================================
% 五层, 自底向上排列
% ============================================================================

% 4 列 x 坐标(相对层 center), 统一用同一套值保证上下列对齐
\newcommand{\colA}{-6.3cm}
\newcommand{\colB}{-2.1cm}
\newcommand{\colC}{2.1cm}
\newcommand{\colD}{6.3cm}

% ---- Layer 1: 仿真场景层 ----
\node[layer] (L1) at (0, 0) {};
\node[layerLabel] at ($(L1.north west) + (0.15cm, -0.25cm)$) {Layer 1 · 仿真场景层 (Webots)};
\node[ours] (L1A) at ($(L1.center) + (\colA, -0.20cm)$) {E-puck 差速模型};
\node[ours] (L1B) at ($(L1.center) + (\colB, -0.20cm)$) {多机场景};
\node[ours] (L1C) at ($(L1.center) + (\colC, -0.20cm)$) {可配置编队};
\node[ours] (L1D) at ($(L1.center) + (\colD, -0.20cm)$) {静态障碍 + 目标};

% ---- Layer 2: 控制接口层 (ROS2) ----
\node[layer, above=0.55cm of L1] (L2) {};
\node[layerLabel] at ($(L2.north west) + (0.15cm, -0.25cm)$) {Layer 2 · 控制接口层 (ROS2)};
\node[upstream] (L2A) at ($(L2.center) + (\colA, -0.20cm)$) {ROS2 Humble};
\node[ours] (L2B) at ($(L2.center) + (\colB, -0.20cm)$) {多机话题契约};
\node[ours] (L2C) at ($(L2.center) + (\colC, -0.20cm)$) {指令分发};
\node[ours] (L2D) at ($(L2.center) + (\colD, -0.20cm)$) {状态广播};

% ---- Layer 3: 算法容器层 (Docker) ----
\node[layer, above=0.55cm of L2] (L3) {};
\node[layerLabel] at ($(L3.north west) + (0.15cm, -0.25cm)$) {Layer 3 · 算法容器层 (Docker)};
\node[upstream] (L3A) at ($(L3.center) + (\colA, -0.20cm)$) {Docker + noVNC};
\node[ours] (L3B) at ($(L3.center) + (\colB, -0.20cm)$) {PyTorch + CUDA};
\node[ours] (L3C) at ($(L3.center) + (\colC, -0.20cm)$) {safepo + safety-gym};
\node[ours] (L3D) at ($(L3.center) + (\colD, -0.20cm)$) {MAPPO 检查点共享};

% ---- Layer 4: 分层集成层 (核心) ----
\node[layer, minimum height=2.6cm, above=0.55cm of L3,
      fill=cBgHighlight, draw=cOursAlt, line width=1pt] (L4) {};
\node[layerLabel, text=black]
  at ($(L4.north west) + (0.15cm, -0.25cm)$)
  {Layer 4 · 分层控制栈集成层 \emph{(本文方法贡献)}};

% 四个集成模块 - 用统一列坐标, 避免重叠
\node[oursAlt] (L4A) at ($(L4.center) + (\colA, -0.35cm)$)
  {MAPPO 策略加载};
\node[oursAlt] (L4B) at ($(L4.center) + (\colB, -0.35cm)$)
  {GCPL 软协调层};
\node[oursAlt] (L4C) at ($(L4.center) + (\colC, -0.35cm)$)
  {CMPL 硬投影层};
\node[oursAlt] (L4D) at ($(L4.center) + (\colD, -0.35cm)$)
  {差速映射};

% Supervisor 容器
\node[draw=cOursAlt, dash pattern=on 2pt off 1pt, rounded corners=3pt,
      fit=(L4A) (L4B) (L4C) (L4D), inner sep=5pt,
      label={[text=black, font=\small\bfseries, fill=white, inner sep=1.5pt, yshift=-3pt]
             below:gcpl\_full\_stack\_supervisor.py}] (L4box) {};

% Layer 4 内部的流向箭头
\draw[flowBold] (L4A.east) -- node[above, font=\footnotesize, text=black, fill=white, inner sep=1pt]{$a^{RL}$} (L4B.west);
\draw[flowBold] (L4B.east) -- node[above, font=\footnotesize, text=black, fill=white, inner sep=1pt]{$a^\star$} (L4C.west);
\draw[flowBold] (L4C.east) -- node[above, font=\footnotesize, text=black, fill=white, inner sep=1pt]{$a_\mathrm{safe}$} (L4D.west);

% ---- Layer 5: 评估管线层 ----
\node[layer, above=0.55cm of L4] (L5) {};
\node[layerLabel] at ($(L5.north west) + (0.15cm, -0.25cm)$) {Layer 5 · 批量评估与可视化管线};
\node[oursEval] (L5A) at ($(L5.center) + (\colA, -0.20cm)$)
  {对比 sweep};
\node[oursEval] (L5B) at ($(L5.center) + (\colB, -0.20cm)$)
  {消融 sweep};
\node[oursEval] (L5C) at ($(L5.center) + (\colC, -0.20cm)$)
  {批量评估};
\node[oursEval] (L5D) at ($(L5.center) + (\colD, -0.20cm)$)
  {多 seed 绘图};

% ============================================================================
% 跨层数据流
% ============================================================================

% ----  跨层主链: 走 layer 框外侧, 用平滑曲线 (Bezier) 贯穿 L1↔L4  ----
%        左侧上行 L1 → L4 (状态), 右侧下行 L4 → L1 (指令)
\draw[flow, rounded corners=8pt]
  ($(L1.west)+(-0.6cm,0)$)
    .. controls ($(L1.west)+(-1.6cm,1.0cm)$) and ($(L4.west)+(-1.6cm,-1.0cm)$) ..
  node[right, pos=0.5, xshift=2pt, font=\footnotesize, text=black, fill=white, inner sep=1pt]{状态上行}
  ($(L4.west)+(-0.6cm,0)$);

\draw[flow, rounded corners=8pt]
  ($(L4.east)+(0.6cm,0)$)
    .. controls ($(L4.east)+(1.6cm,-1.0cm)$) and ($(L1.east)+(1.6cm,1.0cm)$) ..
  node[left, pos=0.5, xshift=-2pt, font=\footnotesize, text=black, fill=white, inner sep=1pt]{指令下行}
  ($(L1.east)+(0.6cm,0)$);

% ----  L3 → L4  (容器提供 runtime, 虚线纵向, 列对齐)  ----
\draw[flow, dashed] (L3A.north) -- (L4A.south);
\draw[flow, dashed] (L3B.north) -- (L4B.south);
\draw[flow, dashed] (L3C.north) -- (L4C.south);
\draw[flow, dashed] (L3D.north) -- (L4D.south);

% ----  L5 → L4  (评估管线触发, 虚线纵向, 列对齐)  ----
\draw[flow, dashed] (L5A.south) -- (L4A.north);
\draw[flow, dashed] (L5B.south) -- (L4B.north);
\draw[flow, dashed] (L5C.south) -- (L4C.north);
\draw[flow, dashed] (L5D.south) -- (L4D.north);

% ============================================================================
% Legend  (放在图下方, 横向排布, 避开五层主体)
% ============================================================================
\node[draw=black!30, rounded corners=2pt, inner sep=5pt, font=\small, text=black,
      anchor=north]
  at ($(L1.south) + (0, -0.4cm)$) (legend)
  {\begin{tabular}{@{}l@{\hspace{3pt}}l@{\hspace{10pt}}l@{\hspace{3pt}}l@{\hspace{10pt}}l@{\hspace{3pt}}l@{\hspace{10pt}}l@{\hspace{3pt}}l@{\hspace{10pt}}l@{\hspace{3pt}}l@{\hspace{10pt}}l@{\hspace{3pt}}l@{\hspace{10pt}}l@{\hspace{3pt}}l@{}}
    \tikz \draw[cOursAlt, line width=1pt, fill=white, rounded corners=1pt]
        (0,0) rectangle (0.4,0.2); & 本文方法贡献 (L4) &
    \tikz \draw[cOurs, line width=1pt, fill=white, rounded corners=1pt]
        (0,0) rectangle (0.4,0.2); & 本文平台扩展 (L1--L3) &
    \tikz \draw[cOursEval, line width=1pt, fill=white, rounded corners=1pt]
        (0,0) rectangle (0.4,0.2); & 评估管线扩展 (L5) &
    \tikz \draw[cUpstream, dashed, fill=white, rounded corners=1pt]
        (0,0) rectangle (0.4,0.2); & 通用基础模块 &
    \tikz \draw[flowBold, shorten >=2pt] (0,0.1) -- (0.4,0.1); & 核心控制管道 &
    \tikz \draw[flow, shorten >=2pt] (0,0.1) -- (0.4,0.1); & 数据/状态流向 &
    \tikz \draw[flow, dashed, shorten >=2pt] (0,0.1) -- (0.4,0.1); & 依赖关系 \\
  \end{tabular}};

\end{tikzpicture}

\end{document}